% ============================================================================
% Chapter 8 — Audio
% ============================================================================
\chapter{Audio}

\textbf{Header:} \texttt{audio/AudioComponents.hpp}\\
\textbf{Namespace:} \texttt{Caffeine::Audio}

The audio system uses ECS components to drive sound playback. Attach an
\texttt{AudioEmitter} to any entity; the \texttt{AudioSystem} picks it up
and plays the referenced clip through SDL3's audio pipeline.

\section{AudioClip}

An \texttt{AudioClip} is a data-only struct that describes loaded PCM audio.
You do not construct these manually — the asset manager fills them when you
load a sound asset. They are referenced by path inside \texttt{AudioEmitter}.

\begin{lstlisting}
struct AudioClip {
    const u8* data          = nullptr;  // raw PCM samples
    u64       size          = 0;        // byte count
    u32       sampleRate    = 44100;
    u16       channels      = 2;        // 1 = mono, 2 = stereo
    u16       bitsPerSample = 16;
    f32       duration      = 0.0f;     // seconds
};
\end{lstlisting}

\section{AudioEmitter}

\begin{lstlisting}
struct AudioEmitter {
    FixedString<128> clipPath;           // path to the .caf audio asset
    f32              volume      = 1.0f; // [0, 1]
    f32              maxDistance = 500.0f;
    bool             loop        = false;
    bool             playOnSpawn = true; // play when entity is created
    bool             spatial     = true; // use positional audio
};
\end{lstlisting}

\begin{longtable}{lp{9cm}}
\toprule
\textbf{Field} & \textbf{Description} \\
\midrule
\texttt{clipPath}    & Path to the audio asset, relative to the asset root. \\
\texttt{volume}      & Master volume for this emitter. \\
\texttt{maxDistance} & Beyond this distance from the listener, volume is 0. Applies only when \texttt{spatial = true}. \\
\texttt{loop}        & If true, restarts automatically when the clip ends. \\
\texttt{playOnSpawn} & If true, playback starts the moment the component is added. \\
\texttt{spatial}     & If true, volume and panning are attenuated by distance to the listener entity. \\
\bottomrule
\end{longtable}

\section{Examples}

\subsection{One-shot sound effect}

\begin{lstlisting}
Entity boom = world.create("Explosion");
world.add<Transform>(boom, Transform{hitPosition});

AudioEmitter sfx;
sfx.clipPath    = "audio/explosion.caf";
sfx.volume      = 0.9f;
sfx.loop        = false;
sfx.playOnSpawn = true;
sfx.spatial     = true;
sfx.maxDistance = 800.0f;
world.add<AudioEmitter>(boom, sfx);
// Destroy after the clip's duration if you want a fire-and-forget entity
\end{lstlisting}

\subsection{Looping background music}

\begin{lstlisting}
Entity music = world.create("BGMusic");
// No transform needed for non-spatial audio
AudioEmitter bgm;
bgm.clipPath    = "audio/soundtrack_01.caf";
bgm.volume      = 0.7f;
bgm.loop        = true;
bgm.spatial     = false;  // 2D: same volume everywhere
bgm.playOnSpawn = true;
world.add<AudioEmitter>(music, bgm);

// Mark as persistent so it survives scene transitions
world.add<PersistentComponent>(music, PersistentComponent{true});
\end{lstlisting}

\subsection{Footstep sounds triggered from a script}

\begin{lstlisting}
// Inside a CppScript (see Chapter 10)
void onUpdate(Entity entity, World& world, f32 dt) override {
    if (!world.has<AudioEmitter>(entity)) {
        AudioEmitter step;
        step.clipPath    = "audio/footstep.caf";
        step.volume      = 0.5f;
        step.loop        = false;
        step.playOnSpawn = false;
        world.add<AudioEmitter>(entity, step);
    }

    if (m_stepCooldown <= 0.0f && isMoving()) {
        auto* emitter = world.get<AudioEmitter>(entity);
        emitter->playOnSpawn = true;   // retrigger
        m_stepCooldown = 0.35f;
    }
    m_stepCooldown -= dt;
}
\end{lstlisting}
